% ============================================================
%  PWr Beamer template – example presentation
%  Compile: pdfLaTeX (Overleaf default)
%  Files: main.tex, pwr-beamer.sty, pwr_logo.png, pwr_partnerzy.png
% ============================================================

\documentclass[aspectratio=169, 11pt]{beamer}
\usepackage{pwr-beamer}

% ---- Metadata ----
% Short title in [...] is shown in the small header line on every slide
\title[Krótki tytuł prezentacji]{Pełny tytuł prezentacji}
\subtitle{Podtytuł (opcjonalnie)}
\course{Nazwa kursu -- Prezentacja 1}          % optional
\author{Imię Nazwisko (nr indeksu)}
\supervisor{dr inż. Imię Nazwisko}              % optional, remove for no supervisor line
\faculty{Wydział Informatyki i Telekomunikacji}
\program{Informatyka Stosowana, studia II stopnia}  % optional
\date{Wrocław, 2026/2027}

% ============================================================
\begin{document}
% ============================================================

% Title slide. For manual line breaks pass the title as an option:
% \pwrtitleframe[Pierwsza linia tytułu\\ druga linia]
\pwrtitleframe

% ------------------------------------------------------------
\begin{frame}{Plan}
  \Large
  \begin{enumerate}\itemsep0.5em
    \item Pierwsza część
    \item Druga część
    \item Podsumowanie
  \end{enumerate}
\end{frame}

% ------------------------------------------------------------
\pwrsection{Pierwsza część}

% ------------------------------------------------------------
\begin{frame}{Slajd z listą}
  \begin{itemize}\itemsep0.4em
    \item Punkt pierwszy z \textbf{wyróżnieniem}
    \item Punkt drugi z \alert{kolorem akcentu}
    \item Punkt trzeci
  \end{itemize}
  \src{Autor, \textit{Tytuł źródła}, rok.}
\end{frame}

% ------------------------------------------------------------
\begin{frame}{Bloki}
  \begin{block}{Blok podstawowy}
    Np. cel pracy albo kluczowa definicja.
  \end{block}
  \vspace{0.2cm}
  \begin{columns}[T,onlytextwidth]
    \begin{column}{0.48\textwidth}
      \begin{exampleblock}{Blok przykładu}
        \begin{itemize}
          \item np. ,,w zakresie''
        \end{itemize}
      \end{exampleblock}
    \end{column}
    \begin{column}{0.48\textwidth}
      \begin{alertblock}{Blok szary}
        \begin{itemize}
          \item np. ,,poza zakresem''
        \end{itemize}
      \end{alertblock}
    \end{column}
  \end{columns}
\end{frame}

% ------------------------------------------------------------
\begin{frame}{Rysunek i tekst w dwóch kolumnach}
  \begin{columns}[c,onlytextwidth]
    \begin{column}{0.47\textwidth}
      \centering
      \begin{tikzpicture}[font=\small]
        \draw[thick, PwrRed, fill=PwrSand!60] (0,0) circle (0.5);
        \draw[-{Stealth}, very thick, PwrDark] (0.6,0) -- (3,0);
        \fill[black!25] (3.2,-0.8) rectangle (3.5,0.8);
      \end{tikzpicture}
      % Or an image: \includegraphics[width=\linewidth]{rysunek.png}
    \end{column}
    \begin{column}{0.5\textwidth}
      \begin{itemize}\itemsep0.4em
        \item[] \textbf{Idea:} opis
        \item[] \textbf{Zaleta:} opis
        \item[] \textbf{Wada:} opis
      \end{itemize}
    \end{column}
  \end{columns}
  % Mark AI-assisted figures/tables (WIT AI guidelines):
  \nsi{Rysunek}
\end{frame}

% ------------------------------------------------------------
\begin{frame}{Tabela}
  \renewcommand{\arraystretch}{1.3}
  \begin{tabular}{@{}lll@{}}
    \toprule
    \textbf{Kolumna 1} & \textbf{Kolumna 2} & \textbf{Kolumna 3} \\
    \midrule
    wartość & wartość & wartość \\
    wartość & wartość & wartość \\
    \bottomrule
  \end{tabular}
\end{frame}

% ------------------------------------------------------------
\begin{frame}{Podsumowanie}
  \begin{itemize}\itemsep0.5em
    \item Najważniejszy wniosek
    \item Drugi wniosek
  \end{itemize}
  \vspace{0.8cm}
  \begin{center}
    {\Large\color{PwrRed}\bfseries Dziękuję za uwagę!}
  \end{center}
\end{frame}

% ------------------------------------------------------------
\begin{frame}{Literatura}
  \footnotesize
  \begin{itemize}\itemsep0.2em
    \item A. Autor, \textit{Tytuł książki}, Wydawnictwo, rok.
    \item B. Autor i in., \textit{Tytuł artykułu}, Czasopismo 1(2), rok.
  \end{itemize}
\end{frame}

% ------------------------------------------------------------
\begin{frame}{Wykorzystanie narzędzi AI}
  \begin{itemize}\itemsep0.4em
    \item \textbf{Narzędzie:} ...
    \item \textbf{Zakres:} ...
    \item \textbf{Oznaczenia:} elementy przygotowane z użyciem NSI opisano w prawym dolnym rogu slajdu
  \end{itemize}
\end{frame}

\end{document}
